% ECONOMICS_PROBLEM: {"id": "E52-316", "title": "Responding to persistent supply shocks", "jel_code": "E52", "source_id": "OP-0248"}
% !TeX program = xelatex
\documentclass[11pt,a4paper]{article}
\usepackage[margin=23mm]{geometry}
\usepackage{fontspec}
\setmainfont{Latin Modern Roman}
\usepackage{amsmath,amssymb,mathtools}
\usepackage{xcolor,enumitem,booktabs,longtable,array,xurl,hyperref}
\definecolor{muted}{HTML}{53636C}
\hypersetup{colorlinks=true,urlcolor=blue,linkcolor=blue}
\setlength{\parindent}{0pt}
\setlength{\parskip}{5pt}
\newcommand{\R}{\mathbb R}
\newcommand{\E}{\mathbb E}
\newcommand{\N}{\mathbb N}
\newcommand{\ind}{\mathbf 1}
\newcommand{\Var}{\operatorname{Var}}
\newcommand{\Cov}{\operatorname{Cov}}
\newcommand{\supp}{\operatorname{supp}}
\DeclareMathOperator*{\argmin}{arg\,min}
\DeclareMathOperator*{\argmax}{arg\,max}
\newcommand{\entrymeta}[3]{{\sffamily\small\color{muted}\textbf{\texttt{#1}}\quad|\quad #2\par #3\par}}
\newcommand{\Problemref}[1]{\texttt{#1}}
\newcommand{\JELref}[2]{\texttt{#2} (JEL \texttt{#1})}
\begin{document}
\section*{Responding to persistent supply shocks}
\textbf{Problem ID:} \texttt{E52-316}\quad \textbf{JEL:} \texttt{E52}\par
% BEGIN ECONOMICS STATEMENT
\label{problem:OP-0248}
{\sffamily\small\color{muted}Original subject group: Inflation and monetary policy\par}
\label{definitions:problem:30007529}
\textbf{Audit classification:} Background; proposed extension. Supply-shock policy guidance is already analyzed; the minimax partial-information problem is editorial. Verification concerns the cited version, not first priority or proof that this exact editorial specification remains unsolved on October 8, 2026.
\subsection*{Definitions and precise research target}
\paragraph{Editorial mathematical specification.} Consider the New Keynesian system
\[
x_t=\mathcal E_t^Px_{t+1}-\sigma^{-1}(i_t-\mathcal E_t^P\pi_{t+1}-r_t^n),\quad
\pi_t=\beta \mathcal E_t^P\pi_{t+1}+\kappa x_t+u_t,\quad
u_{t+1}=\rho u_t+\epsilon_{t+1},
\]
where \(x_t\) is the output gap, \(\pi_t\) inflation relative to target, \(i_t\) the nominal rate, \(r_t^n\) the flexible-price rate, \(\sigma,\kappa>0\), and \(\beta\in(0,1)\). Shock persistence \(\rho\) belongs to a known interval, while the innovation law and possible expectation-anchoring mechanisms belong to a declared compact family \(\mathcal P\). Expectations follow the specified mechanism in each candidate model. The central bank observes noisy indicators and chooses an adapted rule \(r\) satisfying a stated rate bound. Restrict admissible rules to equilibria with finite expected discounted losses.
Characterize the robust rule minimizing
\[
\inf_r\sup_{P\in\mathcal P}E_P\sum_{t\ge0}\beta^t
[\pi_t^2+\lambda x_t^2+\nu(i_t-i_{t-1})^2],
\]
with \(\lambda,\nu\ge0\) fixed policy weights. Determine when learning about persistence justifies accommodation versus stronger stabilization, and how conclusions change when beliefs can become unanchored. Demonstrate computational accuracy and compare robust-policy losses with correctly specified and misspecified benchmarks. This is a model-specific decision problem under explicitly bounded uncertainty; existing qualitative supply-shock guidance does not establish that this precise optimization is an open canonical theorem.

In each candidate model define \(\mathcal E_t^P\) as the specified conditional forecast operator used by private agents; under rational expectations it equals the conditional expectation induced by that equilibrium, while alternative belief models must supply their own kernels. The policymaker's loss is evaluated under the actual law \(P\), not a subjective private forecast law. Restrict \(\rho\) to \([\rho_L,\rho_U]\subset(-1,1)\), give an initial \(u_0,i_{-1}\), and require finite second moments. An admissible policy is measurable with respect to the policymaker's observation filtration and cannot condition on latent shocks. Under model uncertainty the policy must be the same observable-history rule for every candidate. Distinguish commitment at date zero from reoptimization after every observation; they impose different feasible policies and equilibrium incentives.
\subsection*{Variants and assumption changes}
The following are \emph{editorial variants}, not additional published open conjectures.
\begin{enumerate}
\item \emph{Rational-expectations commitment.} Fix \(\rho\), the innovation distribution, and rational expectations; remove parameter uncertainty. Solve the finite-horizon or stabilizing infinite-horizon linear-quadratic benchmark, then quantify the extra loss caused by uncertainty rather than calling the benchmark unsolved.
\item \emph{Learning and anchoring.} Let noisy cost indicators update a prior over \(\rho\), and specify a separate observed anchoring state with its own transition. Compare a Bayesian adapted rule and a minimax rule, keeping the same rate constraints and welfare weights.
\end{enumerate}
\subsection*{References and source audit}
\begin{itemize}
\item Hess Chung; Callum Jones; Antoine Lepetit; Fernando M. Martin. \emph{Implications of Inflation Dynamics for Monetary Policy Strategies}. 2025. Locator: Abstract, printed p.1; Introduction, printed pp.2--3; supply-shock policy analysis. Primary text checked. \url{https://www.federalreserve.gov/econres/feds/files/2025072pap.pdf}.
\end{itemize}
Supply-shock policy guidance is already analyzed; the minimax partial-information problem is editorial.
\begin{enumerate}\raggedright
\item\hypertarget{definitions:source:30007529:1}{} Hess Chung; Callum Jones; Antoine Lepetit; Fernando M. Martin. 2025. \emph{Implications of Inflation Dynamics for Monetary Policy Strategies}. Abstract, printed p.1; Introduction, printed pp.2--3; supply-shock policy analysis.
\url{https://www.federalreserve.gov/econres/feds/files/2025072pap.pdf}.
DOI: \url{https://doi.org/10.17016/FEDS.2025.072}.
\end{enumerate}

\subsection*{Common conventions: Source mathematical conventions}

These conventions apply to each entry unless explicitly replaced.

\textbf{Models and identification.} A primitive model \(m\) specifies all probability laws, feasible choices, information, equilibrium and measurement rules used by its target. The admissible class \(\mathcal M\) is fixed independently of outcomes. If \(P_O(m)\) is its observable probability law and \(\tau(m)\) the target, its identified set at observable law \(P\) is
\[
 \mathcal I_\tau(P)=\{\tau(m):m\in\mathcal M,\ P_O(m)=P\}.
\]
Point identification means that this set is a singleton. An empty set signals incompatibility of data and assumptions. Coordinate bounds need not describe the joint set. Bounds are sharp only if every retained target is attainable or arbitrarily approachable by compatible models. Finite bounds require explicit restrictions on unobserved quantities; unrestricted confounding, hidden wealth, extrapolation or arbitrary equilibrium selection can yield uninformative sets.

\textbf{Causal interventions.} A potential outcome \(Y_i(p)\) is the outcome under a complete policy regime \(p\), including timing, financing, information and market spillovers. The target population and units are fixed before treatment. With interference, use the complete assignment vector or a justified exposure mapping; individual no-interference notation alone cannot identify an economy-wide intervention. A difference of mean potential outcomes does not require the cross-world joint distribution. Conditioning on post-treatment migration, survival, enrollment or employment changes the estimand and needs separate justification.

\textbf{Statistical statements.} An estimator, confidence set, forecast or test specifies a sampling law, model class and loss/coverage criterion. Uniform \((1-\alpha)\)-coverage means \(\inf_{m\in\mathcal M}\Pr_m\{\tau(m)\in C_n\}\geq1-\alpha\), or an explicitly asymptotic version. The whole parameter space trivially has coverage, so informativeness requires a stated width, risk or power objective. A forecasting improvement is relative to a named benchmark, information set and evaluation population.

\textbf{Equilibrium and optimization.} An equilibrium comprises feasible plans, optimality, beliefs where needed, budget constraints and all market/resource clearing. The primitives or a stated benchmark define each correspondence \(\mathcal E(p;m)\). Nonemptiness is assumed or proved, never inferred just from compact policy sets. A selected-equilibrium target uses a declared selection rule; a robust target quantifies over all permitted equilibria. Use suprema or infima unless attainment is established. An \(\varepsilon\)-optimal policy is feasible and within \(\varepsilon\) of the finite optimum value. Report an empty feasible set separately.

\textbf{Welfare and accounting.} Normative utility, interpersonal normalization, social weights, numeraire, discounting and the use of public receipts are primitives. Behavioral utility need not coincide with the welfare criterion. Household welfare includes ownership income and rents once; adding profits again to compensating variation generally double-counts them. A monetary equivalent is defined by an explicit transfer and price convention, with monotonicity and range conditions. Average effects, mechanism contributions, transfers and welfare losses are different targets. Infinite sums require convergence and dynamic feasible plans require terminal or no-Ponzi conditions.

\textbf{Source versions.} A working-paper date, revision date and journal publication year are distinguished. A DOI for a journal version is not automatically the DOI of an earlier manuscript. Locators refer to the stated version and printed pages where specified. A metadata-only check does not verify a claim about a full-text conclusion. Every entry's source subsection narrows the provenance claim accordingly.
% END ECONOMICS STATEMENT
\end{document}
